\documentclass[a4paper,11pt]{article}
\usepackage[margin=26mm]{geometry}
\usepackage{fontspec,amsmath,amsthm,amssymb,booktabs,unicode-math}
\setmainfont{STIXGeneral}\setmathfont{STIX Two Math}
\newtheorem{theorem}{Theorem}\newtheorem{lemma}[theorem]{Lemma}
\newcommand{\Aut}{\operatorname{Aut}}
\title{A corrected construction for unfaithful maps with two flag orbits}
\author{Bounded independent source review}\date{12 October 2026}
\begin{document}\maketitle
\begin{abstract}
We check Propositions 5.9--5.10 of Cunningham, Mochán, Montero, Potočnik
and Toledo, \emph{Faithfulness in maniplexes}, Journal of Algebraic
Combinatorics 64 (2026), article 23. The printed torus construction fails the
maniplex axioms. Replacing its final old color-0 operation by the old color-1
operation repairs it. The converse also needs a choice of the two flag orbits
when a proposed base has repeated adjacent generators. We prove the corrected
statements for possibly infinite flag sets. This reviews physical pages
23--25 and their local background, not the whole article or its final problems.
\end{abstract}

\section{Conventions and the literal defect}
A rank-three maniplex is a connected simple graph with one incident edge of
each color $0,1,2$ at every flag. Its adjacency involutions are fixed-point-free,
have distinct neighbors, and satisfy $r_0r_2=r_2r_0$. Products below act in
written order on the right. An $i$-face is a component after deleting color
$i$. Flags are chain-equivalent when they have the same three faces;
faithfulness means that all these classes are singletons. In class
$2_{\{0,2\}}$ there are precisely two automorphism orbits; colors 0 and 2
preserve them and color 1 interchanges them. These are the source conventions.

On an old finite facet of even length $p=2k$, let $a=r_0,b=r_1$ and
$h=(ab)^k$. The facet is a colored cycle on $2p$ flags, with regular dihedral
action $\langle a,b\rangle$ of order $2p$. Thus $h$ is central. The printed
torus formula makes the new white color 0 equal to $b$ and its color 2 equal
to $ha$. These commute if and only if $a,b$ commute. If $p>2$, distant-color
commutation therefore fails. If $p=2$, then $ha=aba=b$, so colors 0 and 2
have the same white neighbor. Consequently the printed torus construction is
never a simple rank-three maniplex under its stated even-facet hypothesis.

\section{The corrected construction}
Let $K$ have finite even facets, possibly of different lengths and with no
finiteness assumption on its whole flag set. On
$\mathcal F(K)\times\{0,1\}$ define
\[
\begin{array}{lll}
(\Phi,0)R_0=(\Phi a,0),& (\Phi,0)R_1=(\Phi,1),&
(\Phi,0)R_2=(\Phi c,0),\\
(\Phi,1)R_0=(\Phi b,1),& (\Phi,1)R_1=(\Phi,0),&
(\Phi,1)R_2=(\Phi d,1),
\end{array}\tag{1}
\]
where $c=r_2$ and, using $h=(ab)^{p/2}$ on the old facet containing $\Phi$,
\[
 d=h\quad\hbox{for }\widetilde{\operatorname{tr}}_P(K),\qquad
 d=hb\quad\hbox{for }\widetilde{\operatorname{tr}}_T(K).\tag{2}
\]
This changes only the final $r_0$ in the printed torus formula to $r_1$.
The printed projective formula is unchanged.

\begin{theorem}\label{construction}
Both corrected constructions are connected unfaithful rank-three maniplexes.
Digonal facets ($p=2$) are allowed.
\end{theorem}
\begin{proof}
Within an old facet, $h$ and $hb$ are nonidentity involutions commuting
with $b$, and neither equals $b$. For a digon the four dihedral elements
are $1,a,b,ab$, with $h=ab$ and $hb=a$. For longer facets a nonidentity
rotation cannot be a reflection, and $hb=b$ would imply $h=1$. The regular
facet action therefore gives fixed-point-free white matchings distinct from
color 0. On black flags the old color-0 and color-2 matchings already are
distinct and commute. Color 1 changes layer, so coincides with neither.
On white flags $[b,d]=1$. The facet parameter is constant along $a,b$,
so the same checks apply globally when lengths vary.

On the black layer an old $a$-, $b$-, or $c$-step can be followed using
$R_0$, $R_1R_0R_1$, or $R_2$, respectively. Old connectivity and $R_1$
therefore imply new connectivity. On the white layer every word in $a,b$
can be followed using new colors 0 and 1. In particular the distinct
2-adjacent flags $(\Phi,1)$ and $(\Phi d,1)$ share a facet, since
$d\in\langle a,b\rangle$. As they are 2-adjacent, they already share their
vertex and edge. They are chain-equivalent, proving unfaithfulness.
\end{proof}

\section{Regularity and the exact exceptions}
Suppose $K$ is regular. Its flags can be identified with a group
$H=\langle a,b,c\rangle$, where the three marks are distinct nonidentity
involutions, $[a,c]=1$, and $p=|ab|=2k<\infty$. The new flags are
$H\times\{0,1\}$; (1) uses right multipliers, and left multiplication
by $H$ acts transitively on each layer.

\begin{theorem}\label{regularity}
The corrected projective construction is regular exactly when $c=h$;
the corrected torus construction is regular exactly when $c=ha$.
Otherwise its full automorphism group is the displayed $H$ and its type is
$2_{\{0,2\}}$. The exceptions are precisely the source's marked $P\{k\}$
and $T\{k\}$ families where these are simple maniplexes. Thus $P\{1\}$
is allowed, but $T\{1\}$ would have $c=b$ and is excluded by simplicity.
\end{theorem}
\begin{proof}
An automorphism of a connected colored graph is determined by one flag image,
by following finite paths. This holds for infinite graphs too. Hence every
layer-preserving automorphism equals a unique displayed left multiplication.
If any automorphism interchanges the layers, compose with a left multiplication
to obtain $\tau$ sending $(1,0)$ to $(1,1)$. Since it commutes with $R_1$,
it also sends $(1,1)$ to $(1,0)$; together with $H$ it is flag-transitive.

The word $W=(R_0R_1)^p$ multiplies by $h$ on both layers. Indeed,
$(R_0R_1)^2$ multiplies black coordinates by $ab$ and white ones by $ba$;
their $k$th powers coincide. In the projective case $R_2$ and $W$ agree at
$(1,1)$. Commutation with $\tau$ transfers that agreement to $(1,0)$, so
$c=h$. For the torus case use agreement of $R_2$ and $WR_0$ at $(1,1)$;
its transfer gives $c=ha$.

Conversely either equality gives $H=D=\langle a,b\rangle$. The dihedral
automorphism $\theta$ interchanging $a,b$ fixes $h$. In the projective case
it interchanges $c=h$ and $d=h$; in the torus case it interchanges $c=ha$
and $d=hb$. Thus $(g,i)\mapsto(\theta(g),1-i)$ commutes with (1), proving
regularity. The source's defining relations on physical page 19 are exactly
these descriptions of $P\{k\}$ and $T\{k\}$.
\end{proof}

The proof of source Proposition 5.9 labels the displayed quotient
$2_{\{1,2\}}$. Its actual colors give $2_{\{0,2\}}$, as in the statement:
only color 1 changes layer. This is a notation slip, not a further family.

\section{The converse and a degenerate proposed base}
\begin{lemma}\label{dihedral}
Let $a,b$ be distinct involutions in any group. If
$d\in\langle a,b\rangle$ commutes with $b$ and $d\notin\{1,b\}$, then
$p=|ab|$ is finite and even, and
$d\in\{h,hb\}$ where $h=(ab)^{p/2}$.
\end{lemma}
\begin{proof}
Put $t=ab$. Dihedral normal forms are $t^j$ and $t^jb$, with
$btb=t^{-1}$. They do not overlap: if $b\in\langle t\rangle$, the group
would be cyclic and generated by two involutions, forcing $a=b$.
Either form commutes with $b$ exactly when $t^{2j}=1$. If $t$ has infinite
order or odd finite order this gives only $1,b$. With finite even order the
two further possibilities are $h,hb$.
\end{proof}

\begin{theorem}\label{converse}
Every unfaithful rank-three maniplex of full type $2_{\{0,2\}}$ is, up to
duality, one of (1)--(2) from a regular simple rank-three maniplex with finite
even facets. One may need to exchange the names of the two flag orbits
before choosing the base. There is no additional exceptional family.
\end{theorem}
\begin{proof}
We first reprove the rank-three observation used in source Proposition 3.4,
which attributes it to [12, Lemma 3.1]. A chain-equivalent flag lies in the
same edge, whose flags are $\Phi,\Phi^0,\Phi^2,\Phi^{0,2}$. If only the
last is initially known to be equivalent to $\Phi$, it lies in the same
facet, and its 0-neighbor $\Phi^2$ lies there too. Thus $\Phi^2\approx\Phi$.
In every nontrivial case, up to duality, there is a 2-adjacent equivalent pair.

Name its orbit white, choose a root $w$ there and put $v=w^1$ in the black
orbit. Set $H=\Aut(M)$. Each orbit is a free $H$-orbit, since a flag
stabilizer is trivial. Identify $gv$ with $(g,0)$ and $gw$ with $(g,1)$.
There are unique nonidentity involutions $a,c,b,d\in H$ giving the appropriate
0- and 2-neighbors of $v,w$. They are involutions because their squares fix
the base flag. All adjacencies now have (1), and simplicity and distant
commutation give
\[
 a\ne c,\qquad b\ne d,\qquad [a,c]=[b,d]=1.\tag{3}
\]
Connectivity implies $H=\langle a,b,c,d\rangle$: reading a path multiplies
coordinates by these elements, and each is realized by a path ending in its
starting layer (using $R_1$ before and after a white operation if necessary).

The facet through $v,w$ is $D\times\{0,1\}$, where $D=\langle a,b\rangle$.
Since $w^2\approx w$, we have $d\in D$. If $a=b$, then $D=\{1,b\}$,
contradicting (3). Lemma~\ref{dihedral} therefore gives finite even $p=|ab|$
and $d=h$ or $hb$. Consequently $H=\langle a,b,c\rangle$.

If $b\ne c$, the marks $a,b,c$ are pairwise distinct nonidentity involutions
with $[a,c]=1$. Their colored Cayley graph is a regular simple maniplex $K$,
and (1) is its corresponding corrected construction. This uses neither a
finite ambient group nor the intersection property of a string C-group.

The only possible duplicate is $b=c$, omitted in the source's base sentence.
Then $[a,c]=1$ gives $[a,b]=1$, so $p=2$, $H=D\cong C_2^2$ and $h=ab$.
We have $d=ab$ or $d=a$. If $d=a$, interchanging $a,b$ in $C_2^2$ and
interchanging the layers gives a color-preserving automorphism. That would
make $M$ regular, contrary to its hypothesis. Hence $d=ab$.
Exchange the layer names. The four marks become
\[
 (a',b',c',d')=(b,a,ab,b).
\]
Now $a',b',c'$ are the three distinct nonidentity elements of $C_2^2$, and
$h'=a'b'=ab$ satisfies $d'=h'b'$. The base is the colored complete graph
$K_4$ and the construction is the corrected torus version. The new white
pair still shares a facet because $d'\in\langle a',b'\rangle$.
This closes the duplicate-generator case.
\end{proof}

\section{Eight flags refute the printed classification}
Take $H=C_2^2=\{1,x,y,xy\}$ and (1) with
\[
 (a,b,c,d)=(x,y,y,xy).\tag{4}
\]
The layer matchings are distinct and commuting; color 1 changes layer, and
$x,y$ generate all coordinates. Thus this is a simple connected maniplex.
The facet group $\langle x,y\rangle$ and vertex group $\langle y,xy\rangle$
both equal $H$. Each layer is one four-flag edge. There is therefore one
vertex, two edges and one facet, with the layers as chain-equivalence classes.

The word relation $R_1R_0R_1=R_2$ holds at every black flag because $b=c=y$,
and at no white flag because $a=x\ne d=xy$. No color-preserving automorphism
can interchange layers. Left $H$ supplies all four layer-preserving
automorphisms, so the full symmetry type is genuinely $2_{\{0,2\}}$.

If this eight-flag maniplex, or its dual, were a printed construction, its
base would have four flags. Every simple properly 3-colored cubic graph on
four flags is $K_4$, with its three perfect matchings as colors. The printed
torus construction is invalid by Section 1; the projective construction from
$K_4$ is regular by Theorem~\ref{regularity}, since its third mark is $ab=h$.
Neither equals (4) or its dual. Thus literal Proposition 5.10 fails. After the
formula repair and layer choice of Theorem~\ref{converse}, (4) is the
corrected torus construction from $K_4$, not a new family.

The accompanying finite checks are optional diagnostics, not premises for the
infinite statements. No Lean formalization is asserted. The following source
paragraph on exact chain-class sizes, and other sections of F007, remain
outside this bounded review.
\end{document}
